% AutoNeg 1.0 – Kurzhandbuch
% Setzen: pdflatex handbuch.tex (zweimal, für Inhaltsverzeichnis und Verweise)
% Pakete: KOMA-Script, babel, lmodern, booktabs, tabularx, enumitem, csquotes, listings,
% xcolor, hyperref, amssymb, pifont (in jeder üblichen TeX-Distribution enthalten).

\documentclass[a4paper,11pt,parskip=half,DIV=12]{scrartcl}

\usepackage[T1]{fontenc}
\usepackage[utf8]{inputenc}
\usepackage[ngerman]{babel}
\usepackage{lmodern}
\usepackage{textcomp}
\usepackage{amssymb}
\usepackage{pifont}
\usepackage{booktabs}
\usepackage{tabularx}
\usepackage{array}
\usepackage[shortlabels]{enumitem}
\usepackage[autostyle]{csquotes}
\usepackage{xcolor}
\usepackage{listings}
\usepackage[hidelinks]{hyperref}

\setlist{itemsep=0.2em,topsep=0.3em}
\definecolor{hinweis}{RGB}{235,242,250}
\lstset{basicstyle=\ttfamily\small,columns=fullflexible,keepspaces=true,
        backgroundcolor=\color{hinweis},frame=none,xleftmargin=0.5em,
        literate={ä}{{\"a}}1 {ö}{{\"o}}1 {ü}{{\"u}}1 {ß}{{\ss}}1 {–}{{--}}1 {×}{{$\times$}}1}

% Bedienelemente und Tasten
\newcommand{\menu}[1]{\textsf{#1}}                    % Menü, Schaltfläche, Feld
\newcommand{\taste}[1]{\fbox{\footnotesize\textsf{#1}}} % Taste
\newcommand{\datei}[1]{\texttt{#1}}                   % Datei, Ordner, Befehl
\newcommand{\ok}{\ding{51}}
\newcommand{\nein}{\ding{55}}

\title{AutoNeg 1.0}
\subtitle{Kurzhandbuch}
\date{Oktober 2026}

\begin{document}
\maketitle

\noindent AutoNeg wandelt kameradigitalisierte Filme (Farbnegativ C-41, Schwarzweißnegativ und
Dia E-6) aus RAW-Dateien in fertige Positive um. Das Programm erkennt Bildfeld und Filmbasis
automatisch, bietet für Schwarzweiß eine Gradation nach Ilford-Multigrade-Filterstufen,
bearbeitet ganze Rollen und erzeugt einen klassischen Kontaktabzug zum Drucken und Abheften.

Deine Originalscans werden \emph{niemals} verändert, verschoben oder gelöscht. Jede Bearbeitung
liegt in einer kleinen Begleitdatei neben dem Scan und wird sofort gespeichert.

\tableofcontents

% ---------------------------------------------------------------------------------------------
\section{Installation unter Windows}

\paragraph{Voraussetzungen} Windows 10 ab Version 1903 oder Windows 11, 64 Bit. Für RAW-Dateien
mit 60--100 Megapixeln sind 16\,GB Arbeitsspeicher empfehlenswert; mehr Prozessorkerne
beschleunigen Export und Kontaktabzug. Eine Grafikkarte wird in Version 1.0 nicht genutzt.

\paragraph{Installieren}
\begin{enumerate}
  \item \datei{AutoNeg-1.0.0-Setup-x64.exe} starten.
  \item Wählen, ob AutoNeg \menu{für alle Benutzer} (erfordert Administratorrechte) oder
        \menu{nur für mich} installiert wird.
  \item Optional: \menu{Desktop-Symbol erstellen} und \menu{Kommandozeile autoneg-cli zum
        Suchpfad (PATH) hinzufügen}. Die zweite Option braucht man nur, wenn man
        \datei{autoneg-cli} in der Eingabeaufforderung ohne vollständigen Pfad aufrufen möchte
        (Abschnitt~\ref{sec:cli}).
  \item Nach der Installation lässt sich AutoNeg direkt starten. Im Startmenü stehen außerdem
        das Handbuch und die Lizenzen der mitgelieferten Bibliotheken.
\end{enumerate}

\paragraph{Aktualisieren} Eine neue Version einfach über die vorhandene installieren; Einstellungen
und Vorlagen bleiben erhalten.

\paragraph{Deinstallieren} Über \menu{Einstellungen \textrightarrow{} Apps} bzw. die
Systemsteuerung. Einstellungen, Vorlagen (\datei{\%APPDATA\%\textbackslash AutoNeg}) und die
Begleitdateien neben den Scans bleiben dabei bewusst erhalten.

\paragraph{Linux} Ein Installationspaket für Linux folgt mit einer späteren Ausgabe.

% ---------------------------------------------------------------------------------------------
\section{Schnellstart}

\begin{enumerate}
  \item \menu{Datei \textrightarrow{} Ordner öffnen …} (\taste{Strg+O}) und den Ordner mit den
        Scans einer Rolle wählen. Alternativ einen Ordner oder eine Datei auf das Fenster ziehen.
  \item Bei einer neuen Rolle erscheint der \emph{Import-Bildschirm}: Er zeigt die Scans
        unverändert (als Negativ) und fragt, um welchen Film es sich handelt
        (\menu{Farbnegativ}, \menu{S/W-Negativ} oder \menu{Dia}). Filmtyp wählen, optional
        \menu{Bildfeld und Filmbasis gleich erkennen} ankreuzen, dann \menu{Importieren}.
  \item Ohne die Erkennung beim Import: im Filmstreifen alle Bilder markieren (in den
        Filmstreifen klicken, dann \taste{Strg+A}) und \menu{Auto-Bildfeld} klicken. AutoNeg
        erkennt Bildfeld, Schräglage und Filmbasis jedes Bilds; erst mit der Filmbasis wird ein
        Farbnegativ farbrichtig.
  \item Ein Bild anklicken und im Reiter \menu{Anpassen} Belichtung, Farbe und Kontrast
        einstellen. Bei Schwarzweiß im Reiter \menu{Gradation} die Filterstufe wählen.
  \item Gelungene Einstellungen mit \menu{Bearbeiten \textrightarrow{} Bearbeitung kopieren}
        (\taste{Strg+C}) und \menu{Einfügen …} (\taste{Strg+V}) auf weitere Bilder übertragen.
  \item Bilder markieren und \menu{Datei \textrightarrow{} Exportieren …} (\taste{Strg+E}).
  \item Optional: \menu{Datei \textrightarrow{} Kontaktabzug …} (\taste{Strg+K}) für einen
        druckfertigen Kontaktabzug der ganzen Rolle.
\end{enumerate}

% ---------------------------------------------------------------------------------------------
\section{Das Hauptfenster}

Am linken Rand wechselt man zwischen zwei Ansichten:
\begin{description}
  \item[Übersicht] (\taste{G}) alle Bilder der Rolle im aktuellen Stand als Raster, wie ein
        Leuchttisch. Oben steht der \menu{Filmtyp der Rolle}, der sich hier jederzeit ändern
        lässt. Ein Doppelklick öffnet ein Bild in der Bildansicht. Markierungen gelten in beiden
        Ansichten.
  \item[Bildansicht] (\taste{E}) ein Bild bearbeiten, mit diesen Bereichen:
\end{description}

\begin{description}
  \item[Bildfläche] In der Mitte: das aktuelle Bild als Positiv. Bei 100\,\% oder 200\,\%
        verschiebt man den Ausschnitt mit gedrückter linker Maustaste.
  \item[Filmstreifen] Unten: alle Bilder der Rolle mit Bildnummer. \ok{} bedeutet
        \enquote{behalten}, \nein{} \enquote{verworfen}. Mit \taste{Strg} und \taste{Umschalt}
        markiert man mehrere Bilder. Das Auswahlfeld darüber filtert nach Bewertung
        (\menu{Alle Bilder}, \menu{Behalten}, \menu{Verworfen}, \menu{Unbewertet}). Nicht lesbare
        Dateien erscheinen rot.
  \item[Werkzeuge] Rechts, mit den Reitern \menu{Konvertierung}, \menu{Gradation} und
        \menu{Anpassen}.
  \item[Histogramm] Rot, Grün, Blau und Helligkeit. Ein blaues Dreieck oben links zeigt
        beschnittene Tiefen, ein rotes oben rechts beschnittene Lichter.
  \item[Bearbeitungsliste] Alle Änderungen des aktuellen Bilds (Abschnitt~\ref{sec:liste}).
  \item[Statusleiste] Meldungen, bei laufenden Aufträgen ein Fortschrittsbalken und die
        Schaltfläche \menu{Abbrechen}.
\end{description}

Die Werkzeugleiste enthält die häufigsten Befehle: Ordner öffnen, Rückgängig, Wiederholen,
Zuschneiden und Drehen, Staubretusche, Vorher/Nachher, Beschnittwarnung, Exportieren und
Kontaktabzug. Alle Docks lassen sich verschieben, abdocken und über ihr Kontextmenü ein- und
ausblenden.

% ---------------------------------------------------------------------------------------------
\section{Eine Rolle öffnen}

AutoNeg arbeitet mit einem Ordner je Rolle. Gelesen wird nur dieser Ordner, nicht seine
Unterordner; die Bilder erscheinen nach Dateinamen sortiert.

\paragraph{Unterstützte Dateien} RAW-Formate gängiger Kameras (\datei{.arw .srf .sr2 .cr2 .cr3
.nef .nrw .raf .rw2 .orf .pef .dng .rwl .3fr .fff .iiq .x3f}) sowie TIFF und JPEG.

\paragraph{Begleitdateien} Zu jedem Scan legt AutoNeg eine Datei mit der vollen Endung des
Scans plus \datei{.autoneg.json} an, zum Beispiel \datei{DSC01234.ARW.autoneg.json}; Angaben zur
ganzen Rolle stehen in \datei{roll.autoneg.json}. Wer einen Scan umbenennt, verliert nichts:
AutoNeg findet die Bearbeitung über eine Prüfsumme wieder.

\paragraph{Schreibgeschützte Ordner} Kann AutoNeg im Scan-Ordner nicht schreiben (etwa auf einer
DVD oder einem Netzlaufwerk ohne Schreibrecht), legt es die Begleitdateien in seinem
Einstellungsordner ab (Abschnitt~\ref{sec:daten}).

\paragraph{Gleichzeitig geöffnet} Ist eine Rolle bereits in einem anderen AutoNeg-Fenster
geöffnet, wird sie nur lesend angezeigt.

% ---------------------------------------------------------------------------------------------
\section{Konvertierung}

\subsection{Filmtyp}
Der Filmtyp legt fest, wie das Negativ bzw. Dia umgerechnet wird: \menu{Farbnegativ} (C-41),
\menu{S/W-Negativ} oder \menu{Dia} (E-6). Solange keiner gewählt ist, zeigt AutoNeg den Scan
unverändert. Festgelegt wird er beim Import für die ganze Rolle; ändern lässt er sich
\begin{itemize}
  \item für die ganze Rolle in der Übersicht (\menu{Filmtyp der Rolle}) und
  \item für einzelne Bilder im Reiter \menu{Konvertierung}: die Wahl gilt für alle markierten
        Bilder und erscheint als Eintrag \enquote{Filmtyp} in der Bearbeitungsliste.
\end{itemize}
Die Bedienelemente richten sich nach dem Filmtyp: Die \menu{Gradation} nach Ilford Multigrade
gibt es nur bei S/W-Negativ; Weißabgleich und Sättigung sind bei S/W gesperrt, die Filmbasis bei
Dias. Ohne Filmtyp sind Regler und Automatiken gesperrt; Zuschnitt, Drehen und Retusche gehen
immer.

\subsection{Bildfeld und Filmbasis}
\begin{description}
  \item[Auto-Bildfeld] erkennt für alle markierten Bilder das Bildfeld (Zuschnitt), die
        Schräglage und die Filmbasis. Ein vorhandener eigener Zuschnitt bleibt erhalten.
  \item[Auto-Basis] misst die Filmbasis (den unbelichteten Filmrand) automatisch.
  \item[Basis messen] für schwierige Fälle: Schaltfläche klicken und auf dem Bild ein Rechteck
        in den unbelichteten Filmrand ziehen (gelb gestrichelt). Nach dem Loslassen misst
        AutoNeg die Basis in diesem Bereich.
\end{description}
Automatisch ermittelte Werte erscheinen in der Bearbeitungsliste mit der Herkunft
\enquote{automatisch} und lassen sich wie jede andere Änderung entfernen.

\subsection{Zuschneiden und Drehen}
\menu{Ansicht \textrightarrow{} Zuschneiden und Drehen} (\taste{C}) zeigt das ganze,
ungeschnittene Bild mit dem Zuschnittrahmen:
\begin{itemize}
  \item Ecke oder Kante ziehen: Größe ändern.
  \item Innerhalb des Rahmens ziehen: Rahmen verschieben.
  \item Außerhalb des Rahmens ziehen: fein drehen (bis $\pm 45$\textdegree).
\end{itemize}
Beim Loslassen werden Zuschnitt und Feindrehung als ein Schritt übernommen.
Vierteldrehungen und Spiegeln: \menu{90\textdegree{} nach links drehen} (\taste{Strg+[}),
\menu{90\textdegree{} nach rechts drehen} (\taste{Strg+]}), \menu{Waagerecht spiegeln}
(\taste{Strg+H}).

% ---------------------------------------------------------------------------------------------
\section{Gradation für Schwarzweiß}

Der Reiter \menu{Gradation} bildet die Filterstufen von Ilford Multigrade nach:
\begin{itemize}
  \item Die Schaltflächen \menu{00}, \menu{0}, \menu{\textonehalf}, \menu{1} … \menu{5} wählen
        eine Stufe; \menu{ohne Filter} entspricht Stufe~2.
  \item \menu{Fein} stellt die Stufe in Zehntelschritten ein ($-1{,}0$ entspricht 00, bis $5{,}0$).
  \item \menu{Split-Grade}: wie in der Dunkelkammer zwei Belichtungen mit Stufe 00 (weich) und
        Stufe 5 (hart); die beiden Regler geben die Anteile in Prozent an.
\end{itemize}

% ---------------------------------------------------------------------------------------------
\section{Anpassen}

Jede Zeile im Reiter \menu{Anpassen} hat einen Regler und ein Zahlenfeld. Beim Ziehen zeigt die
Bildfläche sofort eine Vorschau; übernommen wird der Wert beim Loslassen, mit dem Mausrad, den
Pfeiltasten oder nach Eingabe im Zahlenfeld. \textbf{Doppelklick auf einen Regler} setzt ihn auf
den neutralen Wert zurück.

\begin{center}
\begin{tabularx}{\textwidth}{@{}lXl@{}}
  \toprule
  Regler & Wirkung & Bereich (neutral) \\
  \midrule
  Belichtung      & Helligkeit in Blendenstufen                 & $-5$ … $+5$\,EV (0) \\
  Farbtemperatur  & Weißabgleich warm/kalt                      & $-300$ … $300$ (0) \\
  Tönung          & Weißabgleich grün/magenta                   & $-300$ … $300$ (0) \\
  Kontrast        & Kontrast um den Mittelton                   & $-1$ … $1$ (0) \\
  Tiefen          & dunkle Bereiche aufhellen                   & $0$ … $1$ (0) \\
  Lichter         & helle Bereiche abdunkeln                    & $0$ … $1$ (0) \\
  Sättigung       & Farbsättigung                               & $-1$ … $1$ (0) \\
  Dynamik         & Sättigung schwacher Farben                  & $-1$ … $1$ (0) \\
  Schärfen        & Unscharfmaskierung, Stärke                  & $0$ … $3$ (0) \\
  Rauschminderung & Stärke der Rauschminderung                  & $0$ … $1$ (0) \\
  \bottomrule
\end{tabularx}
\end{center}

\subsection{Staubretusche}
\menu{Ansicht \textrightarrow{} Staubretusche} (\taste{R}): Jeder Klick auf ein Staubkorn setzt
eine Retuschestelle (cyanfarbener Kreis). AutoNeg füllt sie mit Bildinhalt aus der Nähe. Alle
Stellen eines Bilds bilden einen Eintrag \enquote{Staubretusche} in der Bearbeitungsliste.

\subsection{Ansicht}
\begin{itemize}
  \item \menu{Einpassen} (\taste{Strg+0}), \menu{100\,\%} (\taste{Strg+1}),
        \menu{200\,\%} (\taste{Strg+2}).
  \item \menu{Vorher/Nachher} (\taste{\textbackslash}) zeigt den unkonvertierten Scan.
  \item \menu{Beschnittwarnung} (\taste{J}) markiert Bereiche, in denen Tiefen oder Lichter
        beschnitten werden.
\end{itemize}

% ---------------------------------------------------------------------------------------------
\section{Bearbeitungsliste und Rückgängig}\label{sec:liste}

Die Bearbeitungsliste zeigt jede Änderung des aktuellen Bilds genau einmal, in der Reihenfolge,
in der AutoNeg sie anwendet:
\begin{description}
  \item[Änderung] zum Beispiel Belichtung, Zuschnitt, Gradation.
  \item[Wert] die eingestellten Werte.
  \item[Herkunft] \enquote{manuell}, \enquote{automatisch}, \enquote{kopiert} (mit Angabe des
        Quellbilds) oder \enquote{Rollen-Standard}.
  \item[Aktiv] Häkchen entfernen schaltet die Änderung aus, ohne sie zu löschen (Vergleich).
\end{description}
\menu{Eintrag entfernen} löscht die markierte Änderung mit einem Klick.
\menu{Bearbeiten \textrightarrow{} Bild zurücksetzen} (\taste{Strg+Umschalt+R}) entfernt alle
Änderungen auf einmal. Alles lässt sich mit \menu{Rückgängig} (\taste{Strg+Z}) und
\menu{Wiederholen} (\taste{Strg+Y}) zurücknehmen; der Verlauf gilt bis zum Schließen des
Programms. Gespeichert ist jede Änderung sofort.

% ---------------------------------------------------------------------------------------------
\section{Mehrere Bilder bearbeiten}

Befehle wie Auto-Bildfeld, Einfügen, Bewerten, Vorlagen und Export wirken auf die im
Filmstreifen markierten Bilder, ohne Markierung auf das aktuelle Bild.

\paragraph{Kopieren und Einfügen} \menu{Bearbeitung kopieren} (\taste{Strg+C}) übernimmt die
Liste des aktuellen Bilds. \menu{Einfügen …} (\taste{Strg+V}) öffnet einen Dialog: Dort wählt
man die zu übertragenden Änderungen. Bildbezogene Werte (Zuschnitt, Filmbasis, Retusche) sind
anfangs abgewählt. Mit \menu{Bildfeld und Basis neu erkennen} erkennt AutoNeg diese für jedes
Zielbild neu, statt sie zu übernehmen.

\paragraph{Synchronisieren} Ist \menu{Bearbeiten \textrightarrow{} Synchronisieren}
(\taste{Strg+L}) eingeschaltet, wirkt jede Regler- und Werkzeugänderung auf alle markierten
Bilder gleichzeitig.

\paragraph{Vorlagen} \menu{Als Vorlage speichern …} sichert die aktuelle Bearbeitung ohne
bildbezogene Werte unter einem Namen; \menu{Vorlage anwenden} überträgt sie auf die markierten
Bilder.

\paragraph{Bewerten} \menu{Behalten} (\taste{P}), \menu{Verwerfen} (\taste{X}),
\menu{Bewertung entfernen} (\taste{U}). Der Filmstreifen lässt sich danach filtern; der Export
über die Kommandozeile kann sich auf behaltene Bilder beschränken.

% ---------------------------------------------------------------------------------------------
\section{Exportieren}

Bilder markieren und \menu{Datei \textrightarrow{} Exportieren …} (\taste{Strg+E}) wählen.

\begin{center}
\begin{tabularx}{\textwidth}{@{}lX@{}}
  \toprule
  Feld & Bedeutung \\
  \midrule
  Vorlage     & gespeicherte Exporteinstellungen laden; \menu{Als Vorlage speichern …} sichert
                die aktuellen \\
  Format      & TIFF, JPEG, PNG oder JPEG XL \\
  Bittiefe    & 8 oder 16 Bit (JPEG immer 8 Bit) \\
  Qualität    & 1--100 für JPEG und JPEG XL (Standard 90) \\
  Farbraum    & sRGB, Adobe RGB oder ProPhoto RGB; das Profil wird eingebettet \\
  Lange Kante & Größe der längeren Seite in Pixeln; 0 = Originalgröße \\
  als Graustufen & einkanalige Datei mit Graustufenprofil (für S/W) \\
  Zielordner  & Standard: Unterordner \datei{export} im Rollenordner \\
  Dateiname   & Muster mit Platzhaltern (siehe unten) \\
  \bottomrule
\end{tabularx}
\end{center}

\paragraph{Platzhalter im Dateinamen} \datei{\{roll\}} Rollenname (sonst Ordnername),
\datei{\{frame\}} Bildnummer dreistellig, \datei{\{name\}} Name des Scans ohne Endung,
\datei{\{film\}} Filmtyp. Standard ist \datei{\{roll\}\_\{frame\}\_\{name\}}, also etwa
\datei{Toskana\_007\_DSC01234.tif}.

\paragraph{Im Hintergrund} Der Export läuft im Hintergrund und rechnet mehrere Bilder
gleichzeitig; man kann währenddessen weiterarbeiten. Der Fortschritt steht in der Statusleiste,
\menu{Abbrechen} beendet den Auftrag. Halbe Dateien entstehen dabei nicht.

\paragraph{Vorhandene Dateien} werden nicht überschrieben: Existiert eine Zieldatei bereits,
bricht der Export vor dem ersten geschriebenen Bild mit einer Meldung ab. Dann einen anderen
Zielordner oder ein anderes Namensmuster wählen.

% ---------------------------------------------------------------------------------------------
\section{Kontaktabzug}

\menu{Datei \textrightarrow{} Kontaktabzug …} (\taste{Strg+K}) setzt alle Bilder der Rolle
wie auf einem klassischen Kontaktbogen in Streifen auf eine Seite.

\begin{center}
\begin{tabularx}{\textwidth}{@{}lX@{}}
  \toprule
  Feld & Bedeutung \\
  \midrule
  Vorlage            & gespeicherte Layouts laden; \menu{Als Vorlage speichern} sichert das aktuelle \\
  Filmformat         & Kleinbild 24\texttimes 36, Halbformat 18\texttimes 24, Panorama 24\texttimes 65,
                       Mittelformat 6\texttimes 4,5 / 6\texttimes 6 / 6\texttimes 7 / 6\texttimes 9 \\
  Bilder je Streifen & wie beim Zerschneiden des Films (aus dem Filmformat vorbelegt) \\
  mit Filmrand       & zeigt den Filmrand um jedes Bild \\
  Bildinhalt         & \menu{wie bearbeitet} oder \menu{einheitlich}: gleiche Tonwerte für alle,
                       damit Belichtungsunterschiede sichtbar bleiben \\
  Hintergrund        & schwarz oder weiß \\
  Maßstab            & einpassen oder 1:1 \\
  Papier, quer       & A4, A3 oder Letter, hoch oder quer \\
  Heftrand, Lochmarkierung & Rand zum Abheften (Standard 25\,mm) und Markierungen für den Locher \\
  Rollennummer, Notiz & Text im Kopf des Bogens \\
  \bottomrule
\end{tabularx}
\end{center}

\menu{Als PDF …} speichert ein druckfertiges PDF (Standard \datei{kontaktabzug.pdf} im
Rollenordner), \menu{Drucken …} öffnet den Druckdialog.

% ---------------------------------------------------------------------------------------------
\section{Einstellungen}

\menu{Datei \textrightarrow{} Einstellungen …}:
\begin{description}
  \item[Sprache] Deutsch oder Englisch; wirkt nach einem Neustart.
  \item[Threads] Anzahl der Rechen-Threads; 0 = alle Prozessorkerne (Standard). Wirkt sofort.
  \item[Speicherbudget] Arbeitsspeicher für Zwischenbilder bei Export, Kontaktabzug und
        Miniaturen; Standard 40\,\% des Arbeitsspeichers. Ein kleineres Budget lässt weniger
        Bilder gleichzeitig rechnen, entlastet aber den Rechner.
\end{description}

% ---------------------------------------------------------------------------------------------
\section{Kommandozeile}\label{sec:cli}

\datei{autoneg-cli} verarbeitet eine Rolle ohne Oberfläche, mit denselben Ergebnissen wie das
Programm. Die Bearbeitung kommt aus den Begleitdateien. Aufruf in der Eingabeaufforderung oder
PowerShell. Ohne PATH-Eintrag ruft man es mit vollem Pfad auf, bei einer Installation für alle
Benutzer \datei{\textquotedbl C:\textbackslash Program Files\textbackslash AutoNeg\textbackslash
autoneg-cli.exe\textquotedbl}, nur für den eigenen Benutzer
\datei{\%LOCALAPPDATA\%\textbackslash Programs\textbackslash AutoNeg\textbackslash autoneg-cli.exe}:

\begin{lstlisting}
autoneg-cli info <ordner>
autoneg-cli detect <ordner> [--film color|bw|slide]
autoneg-cli export <ordner> -o <ziel> [--format tiff|jpeg|png|jxl] [--bits 8|16]
            [--quality 1-100] [--profile srgb|adobergb|prophoto] [--long-edge <px>]
            [--gray] [--pattern <muster>] [--preset <exportvorlage>] [--only-kept]
            [--overwrite]
autoneg-cli contact-sheet <ordner> -o <datei.pdf> [--preset <layoutvorlage>]
            [--film-format <name>] [--uniform] [--no-rebate] [--white]
            [--paper a4|a3|letter] [--landscape] [--note <text>]
\end{lstlisting}

\begin{description}
  \item[info] listet die Bilder der Rolle mit Zahl der Änderungen und Bewertung.
  \item[detect] erkennt Bildfeld und Basis für alle Bilder, die noch keinen Zuschnitt haben,
        und speichert das Ergebnis.
  \item[export] wie der Exportdialog. \datei{--preset} lädt eine im Programm gespeicherte
        Exportvorlage, weitere Optionen überschreiben sie. \datei{--only-kept} exportiert nur
        behaltene Bilder, \datei{--overwrite} ersetzt vorhandene Dateien.
  \item[contact-sheet] schreibt den Kontaktabzug als PDF; \datei{--film-format} erwartet den
        Namen wie im Dialog, etwa \datei{\textquotedbl Kleinbild 24×36\textquotedbl}.
\end{description}
\datei{--film color|bw|slide} (Farbnegativ, S/W-Negativ, Dia) gilt für Bilder ohne eigenen
Filmtyp; \datei{detect} trägt ihn in die Begleitdateien ein. Ohne Filmtyp bleibt ein Bild der
unveränderte Scan, und \datei{detect} überspringt es mit einer Meldung.
Für alle Befehle: \datei{--config <ordner>} und \datei{--cache <ordner>} wählen andere
Einstellungs- bzw. Zwischenspeicherordner. \datei{--help} zeigt die Hilfe, \datei{--version}
die Version. Rückgabewert 0 bei Erfolg, 1 bei einem Fehler, 2 bei falschem Aufruf.

\paragraph{Beispiel} Alle behaltenen Bilder einer Rolle als JPEG mit 3000\,Pixeln langer Kante:
\begin{lstlisting}
autoneg-cli export D:\Scans\Toskana -o D:\Export\Toskana --format jpeg --long-edge 3000 --only-kept
\end{lstlisting}

% ---------------------------------------------------------------------------------------------
\section{Wo AutoNeg Daten ablegt}\label{sec:daten}

\begin{description}
  \item[Neben den Scans] \datei{<Scan>.autoneg.json} je Bild und \datei{roll.autoneg.json} je
        Rolle. Diese Dateien gehören zur Rolle: beim Kopieren oder Sichern mitnehmen.
  \item[Einstellungen und Vorlagen] \datei{\%APPDATA\%\textbackslash AutoNeg}: \datei{settings.json},
        Unterordner \datei{presets} (Bearbeitungs-, Export- und Layoutvorlagen) und
        \datei{sidecars} (Begleitdateien schreibgeschützter Ordner).
  \item[Zwischenspeicher] \datei{\%LOCALAPPDATA\%\textbackslash AutoNeg}: Miniaturen und das
        Protokoll \datei{autoneg.log}. Der Ordner darf jederzeit gelöscht werden.
\end{description}
Beschädigte Dateien benennt AutoNeg in \datei{….defekt} um und arbeitet mit Standardwerten
weiter; nichts wird stillschweigend überschrieben.

% ---------------------------------------------------------------------------------------------
\section{Tastenkürzel}

\begin{center}
\begin{tabular}{@{}ll@{\qquad}ll@{}}
  \toprule
  \taste{Strg+O} & Ordner öffnen          & \taste{C}              & Zuschneiden und Drehen \\
  \taste{Strg+E} & Exportieren            & \taste{R}              & Staubretusche \\
  \taste{Strg+K} & Kontaktabzug           & \taste{\textbackslash} & Vorher/Nachher \\
  \taste{Strg+Z} & Rückgängig             & \taste{J}              & Beschnittwarnung \\
  \taste{Strg+Y} & Wiederholen            & \taste{Strg+0}         & Einpassen \\
  \taste{Strg+C} & Bearbeitung kopieren   & \taste{Strg+1}         & 100\,\% \\
  \taste{Strg+V} & Einfügen               & \taste{Strg+2}         & 200\,\% \\
  \taste{Strg+L} & Synchronisieren        & \taste{Strg+[}         & 90\textdegree{} nach links \\
  \taste{Strg+Umschalt+R} & Bild zurücksetzen & \taste{Strg+]}      & 90\textdegree{} nach rechts \\
  \taste{P} & Behalten                    & \taste{Strg+H}         & waagerecht spiegeln \\
  \taste{X} & Verwerfen                   & \taste{U}              & Bewertung entfernen \\
  \taste{G} & Übersicht                   & \taste{E}              & Bildansicht \\
  \bottomrule
\end{tabular}
\end{center}
Auf deutschen Tastaturen erreicht man \textbackslash, [ und ] mit \taste{AltGr}; die Kürzel
folgen dem Tastaturlayout.

% ---------------------------------------------------------------------------------------------
\section{Einschränkungen von Version 1.0}

Folgende Bedienelemente sind vorhanden, haben in dieser Ausgabe aber noch keine Funktion:
\begin{itemize}
  \item \menu{Weißabgleich-Pipette} (Reiter \menu{Anpassen}): den Weißabgleich über
        \menu{Farbtemperatur} und \menu{Tönung} einstellen.
  \item \menu{Messbereich anzeigen} (Reiter \menu{Konvertierung}).
  \item \menu{Auto-Tonwerte}: die Tonwertgrenzen setzt die Konvertierung bereits selbst.
\end{itemize}
Die automatische Bildfelderkennung findet bei manchen Scans kein Bildfeld (Meldung
\enquote{Kein Bildfeld erkennbar}), etwa wenn das Bild bis an die Perforation reicht. Dann den
Zuschnitt mit \menu{Zuschneiden und Drehen} setzen und die Filmbasis mit \menu{Basis messen}
im Filmrand bestimmen; ohne Filmbasis bleibt ein Farbnegativ stark farbstichig.

Außerdem: Der Export der Oberfläche überschreibt keine vorhandenen Dateien (die Kommandozeile
mit \datei{--overwrite} schon); eine Kompression für TIFF und PNG ist nicht wählbar (JPEG XL wird
verlustfrei geschrieben); Filmprofile, automatische Staubentfernung und die Rechnung auf der
Grafikkarte folgen mit Version~1.1.

\end{document}
